\chapter{Literature Review}
\label{chap:literature_review}

This chapter reviews the literature needed to motivate and position the study. Formal estimator definitions are deferred to Chapter~\ref{chap:methodology}.

\section{Monte Carlo pricing and the meaning of efficiency}

\textcite{boyle_1977} established Monte Carlo simulation as a practical option-pricing method. Under risk-neutral pricing, paths of the underlying process are simulated, the discounted payoff is evaluated on each path and the sample mean approximates the option price. The method accommodates path dependence, multiple state variables and payoffs that resist analytical integration \parencite{glasserman_2004}.

For independent observations, variance falls as $1/n$. The convergence rate does not depend explicitly on the dimension of the random input. Dimension nevertheless affects cost: generating a longer path takes more time, payoff variance may change and a learned approximation may become harder to train. A statement about convergence in path count is therefore not equivalent to a statement about computational work.

This point was made early. \textcite{hammersley_morton_1956} argued that efficiency should reflect both sampling variance and the work needed to obtain an estimate. It becomes especially important for learned controls, where training, pilot estimation and network evaluation introduce costs that a variance-reduction ratio does not capture. A complete comparison must ask both how much variance is removed and how much work is required to remove it.

\section{Classical variance reduction}

\subsection{Antithetic variates}

Each Gaussian shock vector $\mathbf{Z}$ is paired with $-\mathbf{Z}$. If the resulting payoffs are negatively correlated, the pair average has lower variance than either payoff alone. The method requires no training and no analytical auxiliary. However, each estimator observation costs two path simulations, so equal-observation and equal-path comparisons answer different questions.

\subsection{Control variates}

A control variate is a quantity computed on the same path as the target payoff whose expectation is known. After centring the control and subtracting it (scaled by a coefficient $\beta$), the adjusted estimator has lower variance provided the control is sufficiently correlated with the payoff. The coefficient is generally unknown and must be estimated from data. \textcite{kim_henderson_2004} study adaptive schemes that tune $\beta$ during simulation. A simpler approach estimates $\beta$ on a separate pilot sample and freezes it before final pricing. This makes the final estimator easier to analyse but exposes the pilot as a setup cost.

\textcite{portier_segers_2019} show that enlarging the space of available controls can improve statistical convergence without improving computational efficiency. Flexibility and cost must be assessed together---a theme that recurs throughout this dissertation.

\section{Geometric Asian control variate}

Under GBM, the geometric average of monitored prices is lognormal, so the geometric Asian option has a closed-form price. \textcite{kemna_vorst_1990} exploit this by using the geometric payoff as a control variate for the arithmetic payoff. Because both averages are computed from the same simulated path, the two payoffs are strongly correlated. The geometric control therefore removes a large fraction of the arithmetic payoff's variance at almost no additional cost.

\textcite{curran_1994} uses the geometric mean differently, conditioning rather than controlling, but the underlying insight is the same: the geometric average carries unusually strong analytical information about the arithmetic average.

This makes GCV a demanding benchmark. An NCV can produce a large gain relative to plain MC and still be unattractive if GCV achieves a larger reduction with negligible setup. At the same time, GCV's strength is specific to this payoff and model. A comparable analytical control may not exist for a different payoff or diffusion. The comparison therefore tests the neural method in a deliberately unfavourable setting: classical structure is unusually strong.

\section{Neural control variates}

A neural network can approximate the pricing integrand, but approximation alone does not yield a valid control variate. The network output must be centred by subtracting its expectation. If that expectation must itself be estimated by Monte Carlo, the original problem has only been displaced.

\subsection{Analytical integration}

\textcite{el_filali_ech-chafiq_lelong_reghai_2021} solve this by restricting the architecture. A single hidden layer of ReLU units with Gaussian inputs produces Gaussian pre-activations whose positive-part expectations are available in closed form. The network expectation follows by linearity. This is the construction used in this dissertation.

Their experiments span arithmetic Asian, basket, digital and worst-of payoffs under Black--Scholes, local-volatility and stochastic-volatility models. They report that learned controls can reduce variance across these problems and include training and pricing time in a combined speed-up measure. This is an important advance over a variance-only comparison, but the aggregated measure does not distinguish one-time training cost from marginal pricing cost. A method with high setup cost and low marginal cost can be inferior for a single valuation yet superior after many valuations; the aggregated figure cannot identify the crossover.

\subsection{Overfitting risk}

\textcite{wan_zhong_xiong_zhu_2020} emphasise that a flexible control can overfit a limited training sample: a low training loss need not imply a low residual variance on new observations. This motivates the strict separation of training, validation and final pricing samples adopted here. The relevant measure is out-of-sample variance, not in-sample approximation error.

\subsection{Architectural trade-off}

The shallow analytical network is fast to integrate but may be less expressive than a deeper network as model complexity grows. A deeper network can approximate the payoff more closely while making exact integration difficult or impossible. This dissertation stays with the shallow architecture to preserve analytical tractability.

\section{Cross-contract reuse}

The economic case for training depends partly on whether the resulting control remains useful when contract parameters change. In the language of \textcite{pan_yang_2010}, frozen NCV reuse is a restricted form of transfer learning: the reference contract supplies the source task and a change in strike, volatility or maturity creates a related target task.

Reuse is not guaranteed. Changing the strike moves the payoff threshold. Changing volatility alters the path distribution. Changing maturity affects discounting, path dynamics and monitoring-date placement. Each change can weaken the relationship between the frozen network output and the target payoff. Re-estimating a target-specific $\beta$ can correct for changes in scale or sign when a useful linear relationship persists, but it cannot recover structure the network never learned. The pilot needed to estimate $\beta$ also creates a target-specific cost.

\textcite{el_filali_ech-chafiq_lelong_reghai_2021} report encouraging robustness, mainly for their dimension-reduction method. Their evidence motivates reuse but does not settle the question addressed here. They do not systematically test frozen reuse of the shallow analytical control across an Asian-option contract grid, nor combine it with an explicit setup-versus-marginal cost model.

The two monitoring profiles ($m=252$ and $m=12$) impose a further constraint. Their networks have different input dimensions and are not directly transferable. Reuse is therefore examined across contracts within a monitoring profile, not between profiles.

\section{Positioning of the present study}

The literature establishes that Monte Carlo is flexible but slow, classical controls can be highly effective when analytical structure is available, and neural controls can learn useful structure provided their expectation remains tractable and performance is measured out of sample.

What remains unclear is how these findings interact when the classical benchmark is unusually strong and neural training is treated as a reusable investment rather than part of a single runtime. This study addresses the gap by comparing MC, AV, GCV and NCV under consistent path accounting; separating setup cost from marginal cost; and testing both contract-specific and frozen controls across two monitoring profiles. The question is not simply whether the network reduces variance, but whether it removes enough residual variation, at low enough marginal cost, to compete with the available classical alternative.
